\chapter{1.6 关于推动服务外包加快转型升级的指导意见(商务部)} % 2020

\begin{center}
  {\small 商服贸发〔2020〕12号}
\end{center}

~\newline

\noindent 各省、自治区、直辖市及计划单列市人民政府，新疆生产建设兵团：

近年来，我国服务外包快速发展，已成为生产性服务出口的主要实现途径，在全球价值链中的地位不断提升。随着新一代信息技术的广泛应用，服务外包呈现出数字化、智能化、高端化、融合化的新趋势。为贯彻落实党中央、国务院关于推进贸易高质量发展的部署要求，推动服务外包加快转型升级，经国务院同意，现提出以下意见。

\section{总体要求}

{ \bf （一）指导思想。 }

以习近平新时代中国特色社会主义思想为指导，全面贯彻党的十九大和十九届二中、三中、四中全会精神，围绕推进贸易高质量发展总体要求，充分发挥服务外包在实施创新驱动和培育贸易新业态新模式中的重要促进作用，加快服务外包向高技术、高附加值、高品质、高效益转型升级，全面提升“中国服务”和“中国制造”品牌影响力和国际竞争力。

{ \bf （二）基本原则。 }

数字引领，创新发展。加强数字技术的开发利用，提高创新能力，加快企业数字化转型，不断向价值链中高端攀升。

跨界融合，协同发展。鼓励服务外包向国民经济各行业深度拓展，加快融合，重塑价值链、产业链和服务链，形成相互渗透、协同发展的产业新生态。

统筹推进，全面发展。服务外包示范城市和具备条件的城市大力发展高技术高附加值业务，鼓励其他城市发展符合地方特色的服务外包业务。

内外联动，协调发展。充分发挥离岸服务外包的技术外溢效应，引进先进技术和模式，大力发展在岸服务外包业务，提升企业竞争力。

{ \bf （三）发展目标。 }

到2025年，我国离岸服务外包作为生产性服务出口主渠道的地位进一步巩固，高技术含量、高附加值的数字化业务占比不断提高，服务外包成为我国引进先进技术提升产业价值链层级的重要渠道，信息技术外包（ITO）企业和知识流程外包（KPO）企业加快向数字服务提供商转型，业务流程外包（BPO)企业专业能力显著增强，服务外包示范城市布局更加优化，发展成为具有全球影响力和竞争力的服务外包接发包中心。

到2035年，我国服务外包从业人员年均产值达到世界领先水平。服务外包示范城市的创新引领作用更加突出。服务外包成为以数字技术为支撑、以高端服务为先导的“服务+”新业态新模式的重要方式，成为推进贸易高质量发展、建设数字中国的重要力量，成为打造“中国服务”和“中国制造”品牌的核心竞争优势。

\section{主要任务}

{ \bf （一）加快数字化转型进程。 }

1.支持信息技术外包发展。将企业开展云计算、基础软件、集成电路设计、区块链等信息技术研发和应用纳入国家科技计划(专项、基金等)支持范围。培育一批信息技术外包和制造业融合发展示范企业。

2.培育新模式新业态。依托5G技术，大力发展众包、云外包、平台分包等新模式。积极推动工业互联网创新与融合应用，培育一批数字化制造外包平台，发展服务型制造等新业态。

3.打造数字服务出口集聚区。依托服务贸易创新发展试点地区和国家服务外包示范城市，建设一批数字服务出口基地。

4.完善统计界定范围。将运用大数据、人工智能、云计算、物联网等新一代信息技术进行发包的新业态新模式纳入服务外包业务统计。

{ \bf （二）推动重点领域发展。 }

5.发展医药研发外包。除禁止入境的以外，综合保税区内企业从境外进口且在区内用于生物医药研发的货物、物品，免于提交许可证，进口的消耗性材料根据实际研发耗用核销。

6.扶持设计外包。建设一批国家级服务设计中心。支持各类领军企业、科研院校开放创新设计中心，提升设计外包能力，支持国家级工业设计中心和国家工业设计研究院开展设计服务外包。实施制造业设计能力提升专项行动。

7.推动会计、法律等领域服务外包。通过双边和区域自贸协定谈判，推动有关国家和地区的会计、法律市场开放。支持会计、法律等事务所为国内企业走出去开展跟随服务，积极研究支持事务所境外发展、对外合作的政策。

8.支持业务运营服务外包。各级政府部门要在确保安全的前提下，不断拓宽购买服务领域。\cl{鼓励企业特别是国有企业依法合规剥离非核心业务，购买供应链、呼叫中心、互联网营销推广、金融后台、采购等运营服务。}

{ \bf （三）构建全球服务网络体系。 }

9.有序增加示范城市。完善服务外包示范城市有进有出的动态管理机制，支持更多符合条件的中西部和东北地区城市创建服务外包示范城市。

10.加大国际市场开拓力度。将更多境外举办的服务外包类展会纳入非商业性境外办展支持范围。支持服务外包企业利用出口信用保险等多种手段开拓国际市场。鼓励向“一带一路”沿线国家和地区市场发包，支持中国技术和标准“走出去”。

11.评估优化出口信贷优惠措施。评估服务外包示范城市服务外包企业承接国际服务外包项目享受的出口信贷优惠措施实施效果，适时向全国推广。

{ \bf （四）加强人才培养。 }

12.大力培养引进中高端人才。推动各省市将服务外包中高端人才纳入相应人才发展计划。鼓励符合条件的服务外包企业对重要技术人员和管理人员实施股权激励。

13.鼓励大学生就业创业。对符合条件的服务外包企业吸纳高校毕业生就业并开展岗前培训的，或为高校毕业生提供就业见习岗位的，按规定给予相应补贴。支持办好相关创新创业大赛，对获奖人员、团队或项目在相关政策方面按规定予以倾斜。

14.深化产教融合。完善包括普通高等院校、职业院校、社会培训机构和企业在内的社会化服务外包人才培养培训体系，鼓励高校与企业开展合作，加快建设新工科，建设一批以新一代信息技术为重点学科的服务外包学院。

{ \bf （五）培育壮大市场主体。 }

15.创新金融支持手段。按市场化原则，充分发挥服务贸易创新发展引导基金作用，带动社会资本加大对服务外包产业投资。支持符合条件的综合服务提供商上市融资。

16.降低企业经营成本。经依法批准，对提高自有工业用地容积率用于自营生产性服务业的工业企业，可按新用途办理相关手续。

17.积极培育国内市场。及时修订《服务外包产业重点发展领域指导目录》，完善重点发展领域的支持政策，引导地方和企业因地制宜发展服务外包业务。

18.大力打造公共服务平台。利用外经贸发展专项资金布局建设一批辐射全国的服务外包公共服务平台。支持服务外包领域智库和行业协会发展，充分发挥其在理论研究和贸易促进中的积极作用。将中国国际服务外包交易博览会办成具有国际影响力的精品展会。

{ \bf （六）推进贸易便利化。 }

19.优化海关监管。逐步将服务外包有关事项纳入国际贸易“单一窗口”。加强对服务外包企业的信用培育，引导更多规范守法的服务外包企业成为海关经认证的经营者（AEO）企业。

20.拓展保税监管范围。在确保有效监管和执行相关税收政策的前提下，研究支持对服务外包示范城市“两头在外”的研发、设计、检测、维修等服务业态所需进口料件试点保税监管。

\section{组织实施}

各部门要从全局和战略高度，深刻认识新时期加快服务外包产业转型升级的重大意义，加强组织领导，完善工作机制，严格组织实施，切实将各项任务落到实处、取得实效。各地区要进一步发挥自主性和积极性，加大对服务外包产业的政策支持力度，促进本地区服务外包产业转型升级，更好地服务于经济转型和开放型经济新体制建设。商务部要加强统筹协调，认真研究解决服务外包产业发展中的新情况新问题，创新工作方法，及时总结经验并推广复制。

~\newline

\rightline{商务部 ~ 发展改革委 ~ 教育部 ~ 工业和信息化部}

\rightline{财政部 ~ 人力资源社会保障部 ~ 海关总署 ~ 税务总局}

\rightline{2020年1月6日}
